% Part: methods
% Chapter: proofs
% Section: resources

\documentclass[../../../include/open-logic-section]{subfiles}

\begin{document}

\olfileid{mth}{prf}{res}

\section{ఇతర వనరులు}

గణితంలో నిరూపణలు ఎలా చేయాలో చెప్పే ఉపయోగకరమైన పుస్తకాలు చాలా ఉన్నాయి. ముఖ్యంగా \emph{How to Read and do Proofs: An Introduction to Mathematical Thought Processes} \citep{Solow2013}, \emph{How to Prove It: A Structured Approach} \citep{Velleman2019} చూడండి. \href{http://www.people.vcu.edu/~rhammack/BookOfProof/BookOfProof.pdf}{\emph{Book of Proof}} \citep{Hammack2013}, \href{https://scholarworks.gvsu.edu/books/7/}{\emph{Mathematical Reasoning}} \citep{Sandstrum2019} అనే నిరూపణ పుస్తకాలు ఆన్‌లైన్‌లో ఉచితంగా అందుబాటులో ఉన్నాయి. తత్వశాస్త్రం చదివేవారికి గణిత తర్కణానికి పరిచయంగా \emph{More Precisely: The Math you need to do Philosophy} \citep{Steinhart2018} ఉపయోగపడవచ్చు.

ఇంటర్నెట్‌లో నిరూపణలపై మరికొన్ని సంక్షిప్త మార్గదర్శకాలు ఉన్నాయి; ఉదా., \href{https://math.berkeley.edu/~hutching/teach/proofs.pdf}{``Introduction to Mathematical Arguments''} \citep{Hutchings2003}, \href{https://eugeniacheng.com/wp-content/uploads/2017/02/cheng-proofguide.pdf}{``How to write proofs''} \citep{Cheng2004}.

\subsection{ప్రేరణనిచ్చే వీడియోలు}

ఇంటిపని చేయడానికి ప్రేరణ లేదా? నిరుత్సాహంగా ఉందా? ఈ వీడియోలు తోడ్పడవచ్చు!

\begin{itemize}
\item \url{https://www.youtube.com/watch?v=ZXsQAXx_ao0}
\item \url{https://www.youtube.com/watch?v=BQ4yd2W50No}
\item \url{https://www.youtube.com/watch?v=StTqXEQ2l-Y}
\end{itemize}

\end{document}
